\documentclass[a4paper,11pt]{article}
\usepackage[utf8]{inputenc}
\usepackage[T1]{fontenc}
\usepackage[french]{babel}
\usepackage[left=2.5cm,right=2.5cm,top=1.8cm,bottom=1.8cm]{geometry}
\usepackage{hyperref}
\usepackage{xcolor}
\usepackage{fontawesome5}
\usepackage{lmodern}

\definecolor{primary}{RGB}{0, 82, 147}
\hypersetup{colorlinks=true, linkcolor=primary, urlcolor=primary}

\begin{document}

\noindent
\textbf{Loïc Aron MBASSI EWOLO} \\
Élève-Ingénieur ENSEA / Polytechnique Yaoundé \\
Cergy, France \\
\faPhone\ (+33) 7 59 52 33 98 \quad \faEnvelope\ \href{mailto:loicaron.mbassiewolo@ensea.fr}{loicaron.mbassiewolo@ensea.fr}

\vspace{0.6cm}

\noindent
\textbf{CEA DAM Île-de-France} \\
À l'attention de M. Gilles KLUTH \\
Bruyères-le-Châtel - 91297 Arpajon

\hfill Cergy, le 23 janvier 2026

\vspace{0.5cm}

\noindent\textbf{Objet : Stage - IA et transfert d'apprentissage pour les plasmas}

\vspace{0.4cm}

Monsieur,

\vspace{0.3cm}

Je suis en deuxième année à l'ENSEA, en double diplôme avec Polytechnique Yaoundé où j'ai suivi une spécialisation en mathématiques appliquées. Votre offre sur le transfert d'apprentissage pour les spectres d'absorption m'intéresse parce qu'elle combine deep learning et rigueur mathématique.

\vspace{0.25cm}

J'ai fait un stage de recherche au Mila (Québec) sur le phénomène de « Grokking » : la généralisation tardive des réseaux de neurones après surapprentissage. J'ai reproduit des expériences sous PyTorch, analysé la convergence mathématiquement et validé les résultats par tests statistiques. Ce travail m'a appris à être rigoureux dans l'entraînement de modèles et à comprendre ce qui se passe réellement dans un réseau de neurones.

\vspace{0.25cm}

Le transfert d'apprentissage, je l'ai pratiqué sur des projets NLP. Pour un modèle de détection de contenus problématiques, j'ai fait du fine-tuning sur BERT et GPT-2, en adaptant des modèles pré-entraînés à un dataset spécifique. J'ai aussi travaillé sur une langue peu dotée (Yemba) où les données de haute qualité sont rares, ce qui m'a confronté aux mêmes problématiques que vous décrivez : comment exploiter des données plus précises quand elles sont limitées.

\vspace{0.25cm}

Sur la partie statistique, ma formation en mathématiques appliquées à Polytechnique Yaoundé m'a donné des bases solides : algèbre linéaire, probabilités, analyse numérique. Je ne connais pas les mélanges de gaussiennes en profondeur, mais je comprends le principe et je suis capable d'apprendre rapidement.

\vspace{0.25cm}

Je n'ai pas de background en physique des plasmas. La spectroscopie et la fusion par confinement inertiel sont des domaines nouveaux pour moi. Cela dit, l'offre mentionne qu'un candidat intéressé par l'IA peut appréhender ces domaines, et c'est exactement ma situation. Ma note de 19.25/20 en physique au bac date, mais elle témoigne d'un intérêt pour les sciences physiques.

\vspace{0.25cm}

La possibilité de poursuivre en thèse m'intéresse. L'apprentissage statistique appliqué à la physique est un domaine où je souhaiterais me spécialiser.

\vspace{0.25cm}

Je suis disponible dès avril 2026 pour 5 mois. Si l'opportunité se présente, je serais également intéressé par un contrat de professionnalisation en dernière année d'école.

\vspace{0.4cm}

Cordialement,

\vspace{0.8cm}

\textbf{Loïc Aron MBASSI EWOLO}

\end{document}
